\begin{env}[5.5.1]
\label{IV.2.s5.5.1-fr}
Recordemos (\textbf 0, 16.3.9) que si $A$, $B$ son dos anillos locales noetherianos,
$k$ el cuerpo residual de $A$, $\varphi:A\to B$ un homomorfismo local, se tiene
\begin{equation}
\dim(B)\leq\dim(A)+\dim(B\otimes_A k).
\tag{5.5.1.1}\label{IV.2.s5.5.1.1-fr}
\end{equation}
\end{env}

De ello se deduce:

\begin{proposition}[5.5.2]
\label{IV.2.s5.5.2-fr}
Sean $X$, $Y$ dos preesquemas localmente noetherianos, $f:X\to Y$ un
morfismo, $x$ un punto de $X$, $y=f(x)$. Entonces se tiene
\begin{equation}
\dim(\mathcal O_x)\leq\dim(\mathcal O_y)
+\dim(\mathcal O_x\otimes_{\mathcal O_y}\boldsymbol{k}(y)).
\tag{5.5.2.1}\label{IV.2.s5.5.2.1-fr}
\end{equation}
\end{proposition}

\oldpage[IV-2]{95}

En particular, si $x$ es un \emph{punto maximal de la fibra} $f^{-1}(y)$, se tiene
\begin{equation}
\dim(\mathcal O_x)\leq\dim(\mathcal O_y)
\tag{5.5.2.2}\label{IV.2.s5.5.2.2-fr}
\end{equation}
puesto que $\mathcal O_x\otimes_{\mathcal O_y}\boldsymbol{k}(y)
=\mathcal O_x/\mathfrak m_y\mathcal O_x$ es el anillo local de $x$ en el
preesquema $f^{-1}(y)$, y es, por tanto, de dimensión $0$ por hipótesis.

Vamos a obtener una fórmula más precisa que (5.5.2.1) cuando se supone que
$f$ es un morfismo \emph{de tipo finito}.

\begin{proposition}[5.5.3]
\label{IV.2.s5.5.3-fr}
Sean $A$ un anillo noetheriano, $T$ una indeterminada, $\mathfrak p$ un ideal
primo de $A$; entonces $\mathfrak p'=\mathfrak pA[T]$ es primo en
$B=A[T]$ y $\mathfrak p'\cap A=\mathfrak p$. Existe una infinidad de
ideales primos de $B$ distintos de $\mathfrak p'$, cuya intersección
con $A$ es $\mathfrak p$; estos ideales son dos a dos no comparables por
inclusión. Además, si $\mathfrak q$ es uno de estos ideales, se tiene
\begin{equation}
\dim(B_{\mathfrak q})=\dim(B_{\mathfrak p'})+1
=\dim(A_{\mathfrak p})+1.
\tag{5.5.3.1}\label{IV.2.s5.5.3.1-fr}
\end{equation}
\end{proposition}

\begin{proof}
Para las primeras afirmaciones, la cuestión se reduce inmediatamente, reemplazando
$A$ por $A/\mathfrak p$ y observando que
$A[T]/\mathfrak pA[T]=(A/\mathfrak p)[T]$, al caso $\mathfrak p=0$; resultan
entonces del hecho de que los ideales primos de $B=A[T]$ cuya intersección
con $A$ se reduce a $0$ son exactamente los que no intersecan la parte
multiplicativa $S=A\setmin\{0\}$ del anillo \emph{íntegro} $A$; ahora bien, se sabe que hay
una biyección creciente del conjunto de estos ideales sobre el conjunto de los
ideales primos de $S^{-1}A[T]=K[T]$, donde $K$ es el cuerpo de fracciones de
$A$ (Bourbaki, \emph{Alg. comm.}, chap. II,
\S~2, n\textsuperscript{o}~5, prop.~11). Además, en virtud de (5.5.1.2), se
tiene
$\dim(B_{\mathfrak q})\leq\dim(A_{\mathfrak p})
+\dim(B_{\mathfrak q}/\mathfrak pB_{\mathfrak q})$,
y si $k$ es el cuerpo de fracciones de $A/\mathfrak p$,
$B_{\mathfrak q}/\mathfrak pB_{\mathfrak q}$ se identifica canónicamente con
$(k[T])_{\mathfrak q}$, que es, por tanto, un anillo de valoración discreta y,
por consiguiente, de dimensión $1$. Finalmente, si
$\mathfrak p=\mathfrak p_0\supset\mathfrak p_1\supset\ldots
\supset\mathfrak p_m$ es una cadena de ideales primos de $A$ de longitud
máxima, los ideales $\mathfrak p_jB$ $(0\leq j\leq m)$ son primos en
$B$, dos a dos distintos y contenidos en $\mathfrak q$; por tanto,
$\dim(B_{\mathfrak q})\geq m+1=\dim(A_{\mathfrak p})+1$ y, en consecuencia,
$\dim(B_{\mathfrak q})=\dim(A_{\mathfrak p})+1$. Esta relación también se
escribe $\operatorname{ht}(\mathfrak q)=\operatorname{ht}(\mathfrak p)+1$;
como $\mathfrak q\ne\mathfrak p'$, se tiene además
$\operatorname{ht}(\mathfrak q)\geq\operatorname{ht}(\mathfrak p')+1
\geq\operatorname{ht}(\mathfrak p)+1$ por definición de la altura de un
ideal primo; esto termina de demostrar (5.5.3.1).
\end{proof}
